\section{Inputs from the fluctuation-scale companion}
\label{sec:companion-interface}

We collect the inputs from \cite{PE}, with their conventions and their
uses in this paper. The moment and comparison estimates are separate inputs
from the exponent theorem. The target temperature is fixed in $(1,\infty)$; the auxiliary common
bases $T_m\to T>1$ are treated below. All replica tests have a fixed finite number of labels, and all moment,
Taylor and recursion orders are fixed before the system size tends to infinity.

\subsection{Covariance and mass normalization}

The Hamiltonian in \eqref{eq:model}, multiplied by $\beta$, has covariance
$mTQ^2/2-T/2$. Adding the independent spin-constant Gaussian of variance
$T/2$ in \eqref{eq:completed-free-energy} gives $mTQ^2/2$, exactly the
covariance of \cite{PE}. The probability spin prior subtracts $m\log2$.
Consequently the centered completed variable in the companion is our
$F_m^c-\mathbf E F_m^c$.

Here is the conversion from our terminal mass $a$ to terminal mass one.
For $u=v/a$, put
\[
 \widetilde b(u)=a^2b(au),\qquad
 \widetilde h(u)=a^2h(au),\qquad
 \widetilde F=aF.
\]
The terminal expression becomes $\log\int e^{aH}$; at an intermediate
mass $v=au$, multiplying the continuation by $a$ changes its power mean
to the one of exponent $u$. Thus, including self corrections,
\[
 f_m^{\rm PE}(\widetilde b,\widetilde h)=a f_m(b,h),\qquad
 \widetilde X_u=X_{au},\quad \widetilde C_u=C_{au},\quad
 H_{au}=a\widetilde H_u,\quad W_{au}=a\widetilde W_u.
\]
In particular $vC_v\preceq H_v\preceq aC_v$. The projection denominator
is unchanged:
\[
 (v/a)\sqrt{m(a^2\Delta h)}=v\sqrt{m\Delta h}.
\]
The same change of variables passes the deterministic-weight pinned
inequality to $dv$ with precisely the right side in \eqref{eq:A1}.
At a bare size endpoint $a=\ell$, $b=T/\ell^2$, the total continuation
is $F_m^c/\ell$, so $x/\ell=m^{-1/6}$ when $m=N e^\theta$.

\subsection{Finite-tree calculus and its probability laws}

The following list specifies both the source statement and the additional
conditions retained when it is used.
\begin{enumerate}
\item \cite[Lemmas~2.1--2.3, equation~(1)]{PE} give the source Hessian,
pressure variation, clock inequality, and single-spin pinned moments in
Lemma~\ref{found:estimates}. The pinned estimate requires zero deterministic
external field, deterministic weights, and the full single-path law. If
other branches are present, their normalized kernels are first marginalized;
remaining random factors are treated afterwards by unconditional H\"older.
The polynomial form \eqref{eq:A2} is stated and proved in
\cite[Appendix, ``Source and pressure calculus; the scope of pinned
moments'']{PE}. Section~\ref{sec:foundations} recalls its proof in our
normalization: for fixed terminal spin the gradient of $P_v(L_v)$ is
$P'_v(L_v)$ times the original gradient. The same block-Hessian
inequality and Brascamp--Lieb argument apply, and gauge symmetry removes
the pin without changing the centering. The coefficients of $P_v$ are
deterministic under the positive path law. This estimate is used in
Lemma~\ref{lem:statistics-centered}.

\item \cite[Lemma~3.5, equation~(13)]{PE} supplies \eqref{eq:A3}.
The anchor is determined by the prefix and branches exterior to the
future being projected. Conditional on those data, that future must
retain its normalized positive kernel. The traversed field length and
its positive lower mass are bounded below by inverse powers of size.
In every application here, the total clocks are at most a fixed power
of $\log m$, as stipulated in Section~\ref{sec:foundations}; hence the
coefficient ratios in the finite projection iteration are polynomial
in the size. The lemma itself does not require the optimized constraints
of PE.
The layer estimates in Sections~\ref{sec:global}--\ref{sec:statistics}
use this version with unit-normalized anchors and then integrate the
anchor magnitude. Their Hessian blocking is the calculation in
\cite[proof of Proposition~3.4, preceding equation~(15)]{PE}.

\item \cite[Proposition~2.5, Lemmas~2.6, 2.8 and~2.9, equation~(2)]{PE}
give the real-power allocation rule, row sums, bounded split densities,
and fixed-order Taylor estimates. Each independent Gaussian increment
covariance remains positive semidefinite throughout each interpolation.
Spin tests and their scalar centers
are held fixed during these interpolations; a varying test requires its
own product-rule derivatives. These are the inputs to
Lemmas~\ref{found:cavity-quotient}--\ref{found:size-increment}, the
coefficient algebra in Proposition~\ref{inv:pair-inverse}, and the
compression identity \eqref{eq:E5}. A row collapses only after its
summand is independent of the new label's detailed attachment.
Allocation order is changed before absolute values are taken. A reused
fork retains its atom; a new fork has a continuous mass variable.

\item \cite[Corollary~2.10]{PE} bounds fixed-order source derivatives.
The terminal-datum derivatives in Lemma~\ref{paths:parameter-regularity} are treated
by its diffusion representation with one bounded smooth insertion;
the source bounds control the remaining fixed derivative slots.
The scalar time identities are those of \cite[Lemma~5.3]{PE}, converted
from time $K/T$ to variance time $K$. They are also used for the
single-coordinate freezing estimate in \eqref{eq:D11}.

\item \cite[Proposition~2.14]{PE} is a limit at each fixed system size,
with fixed finite observed topology. Terminal clocks and every prescribed
cut or revelation side are retained. No rate uniform in system size is
inferred. At a simultaneous jump we choose the joint linear filling
\[
 (b_-,h_-)+t(\Delta b,\Delta h),\qquad 0\le t\le1.
\]
If an intermediate filled revelation is observed, approximations preserve
this filling and its specified observations. Integrated clock identities
refer to that chosen filling; only endpoint laws are independent of it.
Nonnegative contributions of filled jumps can be discarded in trace
upper bounds. Fixed-size limits are used in Lemma~\ref{found:estimates},
the scalar path construction, Lemma~\ref{red:lem-augmentation}, and the
finite-resolution passage in Proposition~\ref{inv:pair-inverse}.
The latter first uses the inverse at a finite resolution with its
uniform absolute attachment bound; it does not require convergence of
the inverse kernels themselves.
\end{enumerate}

\Needspace{4\baselineskip}
For the centered-coupling argument of \cite[Lemmas~2.11--2.13 and~4.5]{PE},
the conditioning sigma-field is generated by the common prefix, the
unaffected exterior branches, and the fixed topology. It contains no
observation of the affected future. This preserves the two separate
conditional marginals. The weighted exterior anchor in
\cite[Lemma~4.8]{PE} averages the exterior joint law rather than
conditioning on its terminal spins. The applications here retain these
conditions. In particular Lemma~\ref{lem:statistics-low-sources} and
Remark~\ref{rem:statistics-earlier-pair} explicitly construct a marginal
law independent of the moving fork before using a centered integral.

\subsection{Optimized comparisons and changed constraints}

The all-mass calls in Sections~\ref{sec:global} and~\ref{sec:budgets}
use \cite[Definition~6.1 and Proposition~6.6]{PE}, with data (3)--(6)
there. Writing $m$ for the site count and $\epsilon_{\rm PE}>0$ for
its fixed power slack, the error parameter obeys
\[
 m^{-2/3+\epsilon_{\rm PE}}d_{\rm PE}\lesssim E
 \le m^{-1/2}d_{\rm PE}^2.
\]
The bare paths and all changed competitors are admissible, total clocks
are uniformly bounded, and a fixed positive matrix margin remains after
the comparison. The upper bound for $d_{\rm PE}^2$ has a fixed positive
gap below the limiting $T$. The endpoint condition is (4) of PE on
every required comparable mass band, with the common base $T_m$ of
\eqref{eq:P9}. The common-base extension is stated and proved in
\cite[Appendix, ``A common prelimit temperature'']{PE}.
Section~\ref{sec:global} verifies its use here and recalls the
substitutions $J_0=T_m\chi^2-1$ in PE (26)--(27) and time $K/T_m$
in PE Lemmas~5.4--5.5. This base stays fixed through each finite
parent/child tree. The macroscopic limit is at the fixed
$T>1$, so \cite[Lemma~5.8 and Proposition~5.9]{PE} apply at that limit.
No uniform assertion near the critical temperature is required.

For clarity, the further source inputs to the modified local optimization
are as follows.
\begin{enumerate}
\item The geometry of \cite[Lemma~3.3, equations~(9)--(10)]{PE} is
rederived for the prescribed density box in Lemma~\ref{bud:geometry}.
The first variation has no state-dependent integrating factor in this
box. The original PE density cap is not assumed for its minimizers.

\item Under the stated outer-regularity hypotheses,
Lemmas~\ref{bud:outer-moments} and~\ref{bud:outer-second} establish
the replacement raw and projection estimates \eqref{eq:C8}--\eqref{eq:C9}
and second moments \eqref{eq:C10}, including the enhanced floor.
Under the additional strong-removal hypotheses, Lemma~\ref{bud:common-star}
gives separate mass and radius estimates and one good calibration cut
for all fixed moment orders. For cap/contact
units these proofs use \cite[Lemmas~3.11--3.12, equations~(20)--(21)]{PE}:
whole-cap trace bounds, controlled mean oscillation, deterministic
weights, and comparable masses. Broad caps are sliced with their
additional trace cost retained. These same inputs are used in
Lemmas~\ref{bud:box-moments}, \ref{bud:free-forks} and
\ref{bud:high-primitive}. The enhanced floor is treated separately by
radius slides; small physical floor mass is never used to discard
radius intervals.

\item The weak-removal argument uses the child-data proof of
\cite[Lemma~6.2]{PE}, the movement inequality of
\cite[Lemma~6.3]{PE}, and feasible-barrier arguments of
\cite[Lemma~6.4]{PE}. The child has the ordinary PE constraints.
The parent movements are proved feasible for the new density box
before using movement. Child success gives agreement and then outer
regularity for the child, as in \cite[Lemma~6.5]{PE}; only then are
\cite[Theorem~5.1 and Proposition~5.9]{PE} used.

\item Proposition~\ref{bud:modified-closure}, used in the strong-removal
argument, reuses the proof of
\cite[Proposition~4.1]{PE}, rather than applying its original
optimization hypotheses to the new density box. The preceding
analytic estimates replace PE (12)--(13), (16)--(19). Fixed-degree
tests remain those of \cite[Definition~4.2]{PE}; bounded row changes,
observed-block truncation, boundary covariance, test-family closure,
and projected replacements are \cite[Lemmas~4.3--4.8]{PE}.
The finite-depth cancellation and calibration are
\cite[Lemmas~4.9--4.10, equations~(23)--(27)]{PE}.
The division is by $\chi K_*\asymp W$, not by a stability gap.
One common good cut is used for the fixed finite recursion and its
required moment orders; these orders are fixed before the size limit.
The resulting moving-target estimate is stated explicitly in
\eqref{eq:bud-coarse-closure}; \eqref{eq:bud-coarse-closure-integrals}
gives its two integrated forms, with the mass and radius normalizations
kept separate. These are the versions of PE (22) used here.
The profile proof is then repeated using \cite[Lemmas~5.3--5.6]{PE}
and PE (28)--(29). Identification of its slope uses
\cite[Lemma~5.8 and Proposition~5.9]{PE} and the macroscopic transfer
in Proposition~\ref{bud:coarse-profile}.
\end{enumerate}

The coarse scalar trial used for \eqref{eq:G3} comes from the bounds
in the proof of \cite[Lemma~5.7]{PE}, with the explicit floor, cap,
and cutoff choices given in Section~\ref{sec:global}. Its proof is
used quantitatively; the limiting variational statement alone would
not provide the needed power saving. The scalar control representation
used for coercivity and pair-kernel positivity is proved in
\cite[proof of Lemma~5.8]{PE} by completing the square in It\^o's
formula.

\subsection{The exponent and the one-dimensional input}

The fixed-temperature exponent theorem \cite[Theorem~1.1]{PE} is
used twice in Section~\ref{sec:reduction}: to initialize the stopped
small-positive-Laplace comparison, and to rule out zero limiting
variance after a power rate of variance convergence has been obtained.
The first use combines the exponent with
\cite[Lemmas~7.6--7.7]{PE}; along the nonbare physical paths the
appropriate exponential augmentation is instead proved in
Lemma~\ref{red:lem-augmentation} from the telescoped pinned density.
That proof uses the full joint positive law and then gauge symmetry.
The one-dimensional tail, positive-tail-mass and small-Laplace bounds
are exactly \cite[Lemma~7.7]{PE}.

The exponent supplies neither a variance prefactor nor a limiting law.
Those conclusions use the summable increment estimate of this paper.
Finally $\operatorname{Var}F_m^c=\operatorname{Var}F_m+T/2$; the common
Gaussian therefore disappears after division by $m^{1/6}$. The theorem
is stated for log partition functions. For the physical total free
energy $-F_m/\beta$, the variance constant is divided by $\beta^2$
and the standardized limiting law is reflected.
